\subsection{Variables}

\subsubsection{Regla: camelCase y nombres descriptivos}

\textbf{Regla:} Toda variable local o parámetro deberá escribirse en \texttt{camelCase} (inicial en minúscula). Queda prohibido el uso de nombres de un solo carácter, salvo el contador de un ciclo simple y no anidado; en ciclos anidados, cada contador deberá tener un nombre descriptivo que indique qué recorre.

\textbf{Descripción:} Microsoft indica expresamente que los parámetros y las variables locales deben usar \texttt{camelCase} y que los nombres de una sola letra deben evitarse salvo en contadores de ciclos simples (Microsoft, 2025c).

\textbf{Código correcto:}
\begin{lstlisting}[style=csharpstyle]
public Tile GetTileAt(int rowIndex, int columnIndex)
{
    for (int row = 0; row < BoardSize; row++)
    {
        for (int column = 0; column < BoardSize; column++)
        {
            if (row == rowIndex && column == columnIndex)
            {
                return _tiles[row, column];
            }
        }
    }

    return null;
}
\end{lstlisting}

\textbf{Código incorrecto:}
\begin{lstlisting}[style=csharpstyle]
public Tile GetTileAt(int rowIndex, int columnIndex)
{
    for (int i = 0; i < BoardSize; i++)
    {
        for (int j = 0; j < BoardSize; j++)
        {
            if (i == rowIndex && j == columnIndex)
            {
                return _tiles[i, j];
            }
        }
    }

    return null;
}
\end{lstlisting}

\subsubsection{Regla: nombres de colecciones en plural}

\textbf{Regla:} Toda variable que almacene una colección de elementos (arreglo, lista, diccionario o cualquier tipo que implemente \texttt{IEnumerable<T>}) deberá nombrarse con un sustantivo en plural que describa los elementos que contiene.

\textbf{Descripción:} Las guías de diseño de .NET recomiendan nombrar las colecciones con un sustantivo plural para distinguirlas a simple vista de una variable que almacena un único elemento (Cwalina \& Abrams, 2008).

\textbf{Código correcto:}
\begin{lstlisting}[style=csharpstyle]
public sealed class GameBoard
{
    private readonly List<Player> players;

    public GameBoard(List<Player> players)
    {
        this.players = players;
    }
}
\end{lstlisting}

\textbf{Código incorrecto:}
\begin{lstlisting}[style=csharpstyle]
public sealed class GameBoard
{
    private readonly List<Player> playerList;

    public GameBoard(List<Player> playerList)
    {
        this.playerList = playerList;
    }
}
\end{lstlisting}

\subsubsection{Regla: una declaración por línea}

\textbf{Regla:} No se declarará más de una variable en la misma sentencia, incluso cuando compartan el mismo tipo de dato.

\textbf{Descripción:} McConnell recomienda declarar cada variable en su propia línea y, de ser necesario, acompañarla de un comentario específico, ya que declarar varias variables juntas dificulta modificarlas o documentarlas de forma independiente (McConnell, 2004).

\textbf{Código correcto:}
\begin{lstlisting}[style=csharpstyle]
int currentRow = startingRow;
int currentColumn = startingColumn;
\end{lstlisting}

\textbf{Código incorrecto:}
\begin{lstlisting}[style=csharpstyle]
int currentRow = startingRow, currentColumn = startingColumn;
\end{lstlisting}

\subsubsection{Regla: prefijos interrogativos para variables booleanas}

\textbf{Regla:} Toda variable de tipo \texttt{bool} deberá nombrarse como una pregunta cerrada, utilizando prefijos como \texttt{is}, \texttt{has}, \texttt{can} o \texttt{should}.

\textbf{Descripción:} Martin sostiene que los nombres booleanos deben leerse como una afirmación o pregunta clara (por ejemplo, \texttt{isActive}) para que el código que los utiliza se lea de forma natural (Martin, 2008).

\textbf{Código correcto:}
\begin{lstlisting}[style=csharpstyle]
bool isPlayerTurnActive = true;
bool hasReachedFinalTile = false;
\end{lstlisting}

\textbf{Código incorrecto:}
\begin{lstlisting}[style=csharpstyle]
bool playerTurnActive = true;
bool finalTileReached = false;
\end{lstlisting}
